\section{Robustness of the study-2 conclusions}
\label{sec:robust}

A second review asked whether the response-derived clusters favour EOST-M, whether the
Bayesian CATs were estimated with enough MCMC, and whether the effects generalize across
benchmarks. Three follow-ups (RF1--RF3), whose plans and decision rules were written before
they were run, answer these questions, together with exploratory extensions (RF2b--RF2d,
RF3b) planned in the same way; none changes a frozen decision.

\paragraph{Clusters that never see the target's responses (RF1).} The strict clusters use
residuals computed from every system's responses, the target's included, and excluding the
calibrators that resemble a target could make its panel harder to calibrate. Keeping the
frozen targets, holdouts and seeds, we rebuilt the panels from heuristic families alone,
from residuals computed on the other benchmarks, and with residual thresholds from 0.4 to 0.7.
Under every rule, R2 and C1 kept their frozen median improvements (\RobRTwoMedMin\% and
\RobCOneMedMin\%), R2 remained positive on all eight benchmarks with \RobRTwoFavRange\
significant, and C1 positive on seven with \RobCOneFavRange\ significant and none against:
the robustness criterion fixed in advance was \RobRuleMet.

\paragraph{MCMC quality of the Bayesian CATs (RF2--RF2d).} We ran the frozen sampler with
four long chains on \ConvPanels\ panels and computed rank-normalized split $\hat R$ and bulk
effective sample sizes \citep{vehtari2021rank}. The overall location converged ($\hat R<1.05$
on \ConvLocRhatFive\% of panels), but $\sigma_b$, $\mu_a$ and $\sigma_a$ did not
($\hat R<1.05$ on at most \ConvHypRhatFiveMax\%): multiplying all abilities and difficulties
by a constant and dividing the discriminations by it leaves the likelihood unchanged, so with
eight calibrators only their $\mathcal N(0,1)$ prior fixes the scale of the calibration, and
the frozen sampler has no move along this ridge. Under the rule fixed in advance the study-2
statement about the Bayesian CATs was \ConvRule: the long chains changed the cell-level log
loss by a median of \ConvLongAbsVsSixtyMin--\ConvLongAbsVsSixtyMax\% relative to a burn-in of 60,
more than the 0.5\% allowed, while the long-chain arms differed from EOST-M by a median of \ConvLongVsEostMin\
to \ConvLongVsEostMax\%, with the benchmarks split. Exploratory samplers then added a
Metropolis--Hastings move along the scale ridge with the item parameters summed out (RF2b, a
partially collapsed Gibbs sampler with a discrete form of parameter expansion;
\citealp{vandyk2008partially,liu1999parameter}) and sampled the model of
Lemma~\ref{lem:exact} exactly, with collapsed moves for the hyperparameters, the scale and the
location (RF2d). RF2d came much closer to convergence (medians of $\hat R$ at most
\ExtRhatMedMaxAll, largest value \ExtRhatMaxAll\ against \ConvHypRhatMax\ for the frozen
sampler) but reached $\hat R<1.05$ on only \ExtRhatFiveMin\% of panels for its slowest
quantity, short of the 90\% fixed in advance, so a fully converged comparator was not
obtained. Across all \AllBayesNArms\ combinations of Bayesian arm and sampler, the median
difference from EOST-M in held-out log loss stayed within $\pm$\AllBayesVsEostMaxAbs\%, with
the benchmarks split: the comparison did not move as the sampler improved, which makes an
artefact of the frozen chain unlikely. The RF2d arms covered the score as often as EOST-M
on the same cells (\ExtCovMin--\ExtCovMax\ against \ExtEostCov), and their score error was lower by a median of
\ExtIaeGainMin--\ExtIaeGainMax\% over cells but on only \ExtIaeBanksMin--\ExtIaeBanksMax\ of
8 benchmarks, a possible advantage that 48 cells cannot establish.

\paragraph{Effects across benchmarks (RF3, RF3b).} The frozen rule requires significance on
five of eight benchmarks after a Holm correction over 88 tests, so an effect can favour
EOST-M everywhere and still be declared not supported. A secondary random-effects
meta-analysis that treats the benchmarks as independent found pooled improvements excluding
zero for \MetaNExclZero\ of the 11 superiority hypotheses, including the
\MetaNExclZeroUnsup\ near misses (\MetaUnsupList). Five benchmarks, however, share the same
395 Open LLM Leaderboard systems. Counting them as one source (RF3b: \MetaSrcK\ sources, with
a bootstrap over the shared systems), \MetaSrcNExclZero\ of the 11 intervals excluded zero
(\MetaSrcExclList), and of the near misses only \MetaSrcUnsupKeptList\ kept an interval above
zero; the dependence through shared systems was small (median correlation \MetaSrcR), and the
loss comes from having four heterogeneous sources. Most effects point the same way across
benchmarks; the analysis does not show that they generalize to a new one.
